% generated from Paper 13/anccft13.tex
\renewcommand{\HS}{\cH_\Sigma}
\chapter[\texorpdfstring{XIII. Universal localization of the degeneracy algebra, the $K_1$-characteristic class, and the
Eisenstein boundary}{XIII. Universal localization of the degeneracy algebra, the K1-characteristic class, and the Eisenstein boundary}]{\texorpdfstring{Algorithmic non-commutative class field theory, XIII: universal localization of the degeneracy algebra, the $K_1$-characteristic class, and the
Eisenstein boundary}{Algorithmic non-commutative class field theory, XIII: universal localization of the degeneracy algebra, the K1-characteristic class, and the Eisenstein boundary}}\label{chap:XIII}
\chaptermark{XIII}
\begin{chapterabstract}
Problem 5.1 of [Ch.~\ref{chap:V}] and its residue [Ch.~\ref{chap:VII}, Rem.~3.1] asked for a single class in $K_1$ of a
localized degeneracy algebra whose evaluations reproduce the snake factorization of [Ch.~\ref{chap:VII}]:
the arc term $q^{\,n_k(q-1)}$ and the Eisenstein index $q$ at every level of the Iwahori tower.
Ore localization is unavailable because the degeneracy algebra is non-commutative and
non-normal ([Ch.~\ref{chap:V}, Prop.~4.1]). This paper constructs the class with Cohn's universal
localization, and finds that the construction reorganizes the whole picture. The degeneracy
order $\cH$ is the $\Zp$-algebra generated by the four degeneracy maps $\ta,\taup,\et,\etp$
acting on the augmentation-zero lattices $\cM^{0}_k=\ker(\deg)\subset\Zp^{n_k}$; the tail
Laplacian factors inside it as $\Delta_k=\ta(\taup-\etp)$, tail norm times degeneracy
coboundary. On $\cM^{0}_k$ the Laplacian is injective with finite cokernel, so
$\Delta=\bigoplus_k\Delta_k$ is $\Sigma$-inverting for $\Sigma=\{\Delta\}$; the universal
localization $\cH\to\HS$ exists, the low-degree localization sequence
$K_1(\cH)\to K_1(\HS)\xrightarrow{\partial}K_0(\cH,\Sigma)\to K_0(\cH)$ holds, and
$[\Theta_{\cH}]:=[\Delta]\in K_1(\HS)$ is the characteristic class. Its content is in two
theorems. The level evaluations $\ev_k\colon K_1(\HS)\to\Qp^{\times}$ give
$\ev_k[\Theta_{\cH}]=\pm\pdet(\Delta_k)=\pm n_k\kappa(Y_k)$, and on $K_4$ these have
$2$-adic order exactly $n_k-3$ at $k=1,2,3$: the $K_1$ class carries the pure arc growth
$\mu p^{k}$ and nothing else. The boundary class $\partial[\Theta_{\cH}]$ at level $k$ is the
cokernel of $\Delta_k$ on $\cM^{0}_k$, which sits in a non-split extension
$0\to\Z/n_k\to\coker\Delta_k^{0}\to\Jac(Y_k)\to0$: the Euler term $-k=-\ord_pn_k+\mathrm{const}$
of [Ch.~\ref{chap:III}] is the Eisenstein subclass $[\Z/n_k]$ of the boundary. Axioms (A1) and (A2) of [Ch.~\ref{chap:V}]
are thus separated inside a single $K$-theoretic object, and (A3) holds on abelian towers,
where $\pdet(\Delta_k)=n\kappa(Y_0)\prod_{\chi\ne1}\det D(\chi)$ is the characteristic
element of [Ch.~\ref{chap:IV}] evaluated at all nontrivial characters (exact on the $(1,1,1)$ tower). What
the specification called the ``non-commutative main conjecture'' $\partial[\Theta]=[X]$ is
the definition of $\partial$; the genuine main conjecture --- an analytic $p$-adic
$L$-function of the tower, defined independently, equal to $[\Theta_{\cH}]$ --- is stated and
left open.
\end{chapterabstract}

\section*{Status of the results}

Every numbered statement is proved. The external inputs are the existence and universal
property of Cohn's universal localization \cite{Co} and the low-degree localization exact
sequence for universal localizations \cite[Ch.~4]{Sch}; the higher sequence of
Neeman--Ranicki \cite{NR} requires stable flatness, which we do not verify and do not use.
Everything else is finite linear algebra over $\Zp$, machine-verified exactly on the $K_4$
Iwahori tower at $k\le3$ and on the $(1,1,1)$ abelian tower (Table~\ref{XIII:tab:battery}).
Section~\ref{XIII:sec:scope} says what the main conjecture is and is not.

\section{The degeneracy order and the augmentation-zero lattices}\label{XIII:sec:order}

Setting as in [Ch.~\ref{chap:V}], [Ch.~\ref{chap:VII}]: $Y_k$ the Iwahori path tower over a $(q{+}1)$-regular base,
$q=p$, $n_k=n(q{+}1)q^{k-1}$, $A_k$ the level Hecke operator, $\Delta_k=qI-A_k$,
$\Tor_k=\Jac(Y_k)=\Tor\coker\Delta_k$; $\ta,\et\in M_{n_k\times n_{k+1}}(\Z)$ the tail and head
pushforwards, $\taup,\etp$ their transposes, with [Ch.~\ref{chap:V}, Prop.~1.2]
\begin{equation}\label{XIII:eq:calc}
\ta\etp=A_k,\qquad \etp\ta=A_{k+1},\qquad \ta\taup=\et\etp=qI_{n_k},
\end{equation}
and $\deg\colon\Z^{n_k}\to\Z$ the coordinate sum.

\begin{lemma}[Degree compatibility and the Laplacian factorization]\label{XIII:lem:deg}
$\deg\circ\ta=\deg\circ\et=\deg$ and $\deg\circ\taup=\deg\circ\etp=q\deg$; hence all four
degeneracy maps preserve the augmentation-zero lattices $\cM^{0}_k:=\ker(\deg)\subset\Zp^{n_k}$.
Moreover
\[
\Delta_k=\ta\,(\taup-\etp)=\ta\,\partial_k^{*},
\]
the tail norm composed with the degeneracy coboundary $\partial_k^{*}$ of {\rm[Ch.~\ref{chap:V}, Thm.~2.1]};
in particular $\Delta_k$ lies in the algebra generated by the degeneracy maps and maps
$\Zp^{n_k}$ into $\cM^{0}_k$.
\end{lemma}

\begin{proof}
Pushforwards sum over fibres and pullbacks have $q$ elements per fibre. By \eqref{XIII:eq:calc},
$\ta\taup-\ta\etp=qI-A_k$. Column sums of $\Delta_k$ vanish (in-degree $q$). Verified exactly
at $k=1,2,3$ (Table~\ref{XIII:tab:battery}, lines (I),(B)).
\end{proof}

\begin{definition}[Degeneracy order]\label{XIII:def:order}
For $K\ge1$ let $\cM^{0}_{\le K}:=\bigoplus_{k\le K+1}\cM^{0}_k$, and let
$\cH_K\subset\End_{\Zp}(\cM^{0}_{\le K})$ be the $\Zp$-subalgebra generated by the degeneracy
maps between consecutive levels $\le K{+}1$, restricted to the augmentation-zero lattices by
Lemma~\ref{XIII:lem:deg}. It is an order in a semisimple $\Qp$-algebra (a $\Zp$-subalgebra of a
matrix algebra, finitely generated as a module), non-commutative and non-normal since
$A_kA_k^{t}\ne A_k^{t}A_k$ inside it {\rm[Ch.~\ref{chap:V}, Prop.~4.1]}. Put
\[
\Delta_{(K)}:=\bigoplus_{k\le K}\Delta_k^{0}\ \oplus\ q\cdot\mathrm{id}_{\cM^{0}_{K+1}}\ \in\cH_K,
\qquad \Delta_k^{0}:=\Delta_k|_{\cM^{0}_k},
\]
where the last summand $q\,\mathrm{id}=\ta\taup$ keeps the element inside $\cH_K$. We write
$\cH$ for $\cH_K$ with $K$ fixed and large, and $\Sigma:=\{\Delta_{(K)}\}$.
\end{definition}

\begin{proposition}[$\Sigma$-inverting]\label{XIII:prop:sigma}
$\Delta_k^{0}$ is injective with finite cokernel and invertible over $\Qp$, with
\[
\det\Delta_k^{0}=\pm\,\pdet(\Delta_k)=\pm\,n_k\kappa(Y_k),
\]
the product of the nonzero eigenvalues of $\Delta_k$. Hence $\Delta_{(K)}\otimes\Qp$ is
invertible while $\Delta_{(K)}$ is not invertible in $\cH$, its determinant not being a
$\Zp$-unit: localization is needed and is non-trivial.
\end{proposition}

\begin{proof}
$\ker\Delta_k=\Zp\mathbf1$ and $\deg\mathbf1=n_k\ne0$, so $\Delta_k$ is injective on
$\ker(\deg)$; $\Delta_k$ maps $\Zp^{n_k}$ into $\ker(\deg)$, so $\Delta_k^{0}$ is an
endomorphism of the rank-$(n_k-1)$ lattice $\cM^{0}_k$ with finite cokernel. Its determinant
is the product of the eigenvalues of $\Delta_k$ on the complement of the kernel, i.e.\ the
pseudo-determinant, which equals $n_k\kappa(Y_k)$ by [Ch.~\ref{chap:III}, eq.~(2.2)]. On $K_4$:
$4608=12\cdot384$, $18874368=24\cdot786432$, $316659348799488=48\cdot6597069766656$, with
$2$-adic orders $9,21,45$.
\end{proof}

\section{Universal localization and the characteristic class}\label{XIII:sec:loc}

\begin{definition}[Universal localization and the relative group]\label{XIII:def:loc}
By Cohn \cite{Co} there is a ring homomorphism $\lambda\colon\cH\to\HS$, universal among ring
homomorphisms under which $\Delta_{(K)}$ becomes invertible. Let $K_0(\cH,\Sigma)$ be the
relative group of triples $(P,Q,f)$, $P,Q$ finitely generated projective $\cH$-modules and
$f\colon P\to Q$ a map with $\lambda(f)$ invertible, modulo the usual relations
\cite[Ch.~4]{Sch}. Then
\begin{equation}\label{XIII:eq:seq}
K_1(\cH)\longrightarrow K_1(\HS)\xrightarrow{\ \partial\ }K_0(\cH,\Sigma)\longrightarrow K_0(\cH)
\end{equation}
is exact \cite[Ch.~4]{Sch}, and $\partial$ sends the class of $\lambda(f)$, for
$f\colon P\to Q$ with $\lambda(f)$ invertible, to $[(P,Q,f)]$. Forgetting to modules,
$[(P,Q,f)]\mapsto[\coker f]-[\ker f]$ when these are finite.
\end{definition}

\begin{definition}[Characteristic class]\label{XIII:def:theta}
$[\Theta_{\cH}]:=[\lambda(\Delta_{(K)})]\in K_1(\HS)$. It is well defined by
Proposition~\ref{XIII:prop:sigma}, and $\HS\ne0$ because $\lambda$ extends to the representation
$\cH\to\End_{\Qp}(\cM^{0}_{\le K}\otimes\Qp)$ in which $\Delta_{(K)}$ is invertible.
\end{definition}

\begin{theorem}[The tautological identity, and what it says]\label{XIII:thm:tauto}
$\partial[\Theta_{\cH}]=[(\cM^{0}_{\le K},\cM^{0}_{\le K},\Delta_{(K)})]\in K_0(\cH,\Sigma)$, whose
image under the forgetful map is
\[
\sum_{k\le K}\bigl[\coker\Delta_k^{0}\bigr]+\bigl[\cM^{0}_{K+1}/q\cM^{0}_{K+1}\bigr].
\]
This is the definition of $\partial$; it is the statement the specification called the
``non-commutative main conjecture'', and it has no content beyond
Definition~\ref{XIII:def:loc}. The content is in Theorems~\ref{XIII:thm:eval} and \ref{XIII:thm:boundary}.
\end{theorem}

\section{Evaluation and the Eisenstein boundary}\label{XIII:sec:eval}

For each $k\le K$ the projection $\cH\to\End_{\Zp}(\cM^{0}_k)$ followed by
$\otimes\Qp$ is a ring homomorphism $\cH\to M_{n_k-1}(\Qp)$ under which $\Delta_{(K)}$ becomes
invertible; by universality it factors through $\HS$, and composing with the determinant
gives the \emph{level evaluation}
\[
\ev_k\colon K_1(\HS)\longrightarrow K_1\bigl(M_{n_k-1}(\Qp)\bigr)=\Qp^{\times}.
\]

\begin{theorem}[Level evaluations: pure arc growth]\label{XIII:thm:eval}
$\ev_k[\Theta_{\cH}]=\det\Delta_k^{0}=\pm\,n_k\kappa(Y_k)$, and for all $k\ge1$
\[
\frac{\ev_{k+1}[\Theta_{\cH}]}{\ev_k[\Theta_{\cH}]}=\pm\,q^{\,n_k(q-1)},\qquad
\ord_p\ev_k[\Theta_{\cH}]=n_k+c=\frac{n(q{+}1)}{q}\,p^{k}+c ,
\]
with a constant $c$ depending only on the base. The arc term $\mathrm{Det_{arc}}=q^{\,n_k(q-1)}$
of {\rm[Ch.~\ref{chap:VII}, Cor.~2.2]} is the ratio of consecutive evaluations of the single class
$[\Theta_{\cH}]$; the Euler term does not appear. On $K_4$: $c=-3$ and
$\ord_2\ev_k=9,21,45$ at $k=1,2,3$.
\end{theorem}

\begin{proof}
The first equality is Proposition~\ref{XIII:prop:sigma}. The ratio: $\pdet(\Delta_{k+1})/\pdet(\Delta_k)
=q^{\,n_{k+1}-n_k}=q^{\,n_k(q-1)}$ because the nonzero spectrum of $A_{k+1}$ is that of $A_k$
plus $n_{k+1}-n_k$ zeros [Ch.~\ref{chap:III}, eq.~(2.1)], each contributing the eigenvalue $q$ of
$\Delta_{k+1}$. Equivalently this is the residue constancy $n_k\kappa_kq^{-n_k}=\mathrm{const}$
of [Ch.~\ref{chap:IX}, Thm.~4.1(iii)]. Summing, $\ord_p\ev_k=n_k+c$ with $n_k=n(q{+}1)q^{k-1}$.
\end{proof}

\begin{theorem}[The Eisenstein boundary]\label{XIII:thm:boundary}
For every $k\ge1$ there is an exact sequence of finite $\Zp$-modules
\[
0\longrightarrow\Z/n_k\longrightarrow\coker\Delta_k^{0}\longrightarrow\Jac(Y_k)\longrightarrow0,
\]
so that in the Grothendieck group of finite $\Zp$-modules the level-$k$ component of
$\partial[\Theta_{\cH}]$ is
\[
\bigl[\coker\Delta_k^{0}\bigr]=\bigl[\Jac(Y_k)\bigr]+\bigl[\Z/n_k\bigr],
\qquad
\ord_p|\Jac(Y_k)|=\ord_p\ev_k[\Theta_{\cH}]-\ord_pn_k=\frac{n(q{+}1)}{q}p^{k}-k+\nu_0 .
\]
The Euler term $-k$ of {\rm[Ch.~\ref{chap:III}, Thm.~2.2]} is the Eisenstein subclass $[\Z/n_k]$ of the
boundary, of $p$-adic order $\ord_pn_k=k+\ord_p(n(q{+}1))-1$. On $K_4$ the extension is
non-split at $k=1,2,3$: $\coker\Delta_k^{0}=\Z/8\oplus(\Z/24)^{2}$,
$(\Z/2)^{9}\oplus\Z/16\oplus(\Z/48)^{2}$, $(\Z/2)^{12}\oplus(\Z/4)^{9}\oplus\Z/32\oplus(\Z/96)^{2}$,
against $\Jac(Y_k)\oplus\Z/n_k$, whose invariant factors differ (at $k=1$:
$\Z/2\oplus\Z/4\oplus(\Z/24)^{2}$).
\end{theorem}

\begin{proof}
$\Delta_k(\Zp^{n_k})\subseteq\ker(\deg)$, so $\Jac(Y_k)=\Tor\coker\Delta_k=\ker(\deg)/\Delta_k
(\Zp^{n_k})$ by [Ch.~\ref{chap:VII}, Lem.~1.1] and [Ch.~\ref{chap:IV}, Lem.~1.3]. The lattice $\Delta_k(\Zp^{n_k})$ contains
$\Delta_k(\ker\deg)$ with quotient $\Delta_k(\Zp^{n_k})/\Delta_k(\ker\deg)\cong\Zp^{n_k}/
(\ker\deg+\Zp\mathbf1)\cong\Z/\deg(\mathbf1)=\Z/n_k$, since $\Delta_k$ is injective modulo
$\Zp\mathbf1$. The exact sequence follows, and orders multiply: $|\coker\Delta_k^{0}|=n_k
|\Jac(Y_k)|$, consistent with Proposition~\ref{XIII:prop:sigma}. The $p$-adic law is
Theorem~\ref{XIII:thm:eval} minus $\ord_pn_k$; on $K_4$, $n_k-3-(k+1)=6\cdot2^{k}-k-4$, the law of
[Ch.~\ref{chap:III}, Thm.~2.2]. Non-splitting: compare invariant factors (Table~\ref{XIII:tab:battery}, line (B)).
\end{proof}

\begin{remark}[Axioms (A1), (A2), (A3) in $K$-theory]\label{XIII:rem:axioms}
[Ch.~\ref{chap:V}, \S5] asked for a characteristic formalism producing an arc-density term $\mu$ (A1), an
Euler term $-\chi k$ intrinsic to the base (A2), and the commutative characteristic ideal on
abelian sub-towers (A3); [Ch.~\ref{chap:VII}] delivered all three in evaluated form and left the lift to a
single $K_1$ class as the residue. Theorems~\ref{XIII:thm:eval}--\ref{XIII:thm:boundary} are that lift:
the class $[\Theta_{\cH}]$ evaluates to the pure arc law (A1); its boundary contains the
Eisenstein subclass $[\Z/n_k]$ whose order is the vertex count, and $-\ord_pn_k$ is the Euler
term (A2) --- the fourth reading of $-\chi k$ after [Ch.~\ref{chap:III}], [Ch.~\ref{chap:VII}] and [Ch.~\ref{chap:IX}], and the first in
which it is a class rather than a number; (A3) is Theorem~\ref{XIII:thm:abelian}. The three
axioms are properties of one object, not three formulas.
\end{remark}

\begin{theorem}[Abelian reduction]\label{XIII:thm:abelian}
For an abelian $\Z/\ell^{k}$ voltage tower of {\rm[Ch.~\ref{chap:IV}]} with pencil $D(t)$ and
$\Delta_k=D(S_{\ell^{k}})$ {\rm[Ch.~\ref{chap:IV}, Lem.~1.1]}, the same construction (the group ring
$\Zl[G_k]$ playing the role of $\cH$) gives
\[
\ev_k[\Theta]=\pdet(\Delta_k)=n_k\kappa(Y_k)=n\,\kappa(Y_0)\prod_{\chi\ne1}\det D(\chi)
=n\,\kappa(Y_0)\prod_{\zeta^{\ell^{k}}=1,\ \zeta\ne1}\zeta^{-d}g(\zeta-1),
\]
the characteristic element $g=\det D(1+T)$ of {\rm[Ch.~\ref{chap:IV}, Thm.~2.1]} evaluated at all
nontrivial characters, times the Eisenstein constant $n\kappa(Y_0)$; the Eisenstein
subclass is $\Z/n_k=\Z/n\ell^{k}$ and $\ord_\ell n_k=k+\ord_\ell n$ is the ``$-k$'' of
{\rm[Ch.~\ref{chap:II}, Rem.~5.4]}. Exact on the $(1,1,1)$ tower, $\ell=2$, $k\le3$: $3072$, $2408448$,
$362577395712$ (Table~\ref{XIII:tab:battery}, line (A3)).
\end{theorem}

\begin{proof}
The block-character decomposition [Ch.~\ref{chap:IV}, Lem.~1.1] splits the nonzero spectrum of
$\Delta_k$ into the spectra of $D(\chi)$, $\chi\ne1$, and the nonzero spectrum of
$D(1)=\Delta_0$, whose product is $n\kappa(Y_0)$ [Ch.~\ref{chap:III}, eq.~(2.2)]; $\det D(\zeta)=\zeta^{-d}
g(\zeta-1)$ as in [Ch.~\ref{chap:IX}, Lem.~2.1].
\end{proof}

\section{The verification battery}\label{XIII:sec:battery}

\begin{table}[ht]
\centering\scriptsize
\renewcommand{\arraystretch}{1.2}
\resizebox{\textwidth}{!}{%
\begin{tabular}{lcl}
\toprule
check & levels & result\\
\midrule
(I) \eqref{XIII:eq:calc}; $A_kA_k^{t}\ne A_k^{t}A_k$; degree compatibility of $\ta,\et,\taup,\etp$ & $k=1,2,3$ & exact\\
(S) $\Delta_k^{0}$ injective, invertible over $\Q$, $\det$ not a $\Z_2$-unit ($2$-adic orders $9,21,45$) & $k=1,2,3$ & exact\\
(E) $\det\Delta_k^{0}=n_k\kappa_k=\pdet$; ratios $2^{12},2^{24}$; $\ord_2=n_k-3$ & $k=1,2,3$ & exact (pdet numeric)\\
(B) $|\coker\Delta_k^{0}|/|\Jac(Y_k)|=n_k$; structures; extension non-split; $\Delta_k=\ta(\taup-\etp)$ & $k=1,2,3$ & exact\\
(A3) $(1,1,1)$, $\ell=2$: $\pdet(\Delta_k)=n\kappa_0\prod_{\chi\ne1}\det D(\chi)$ & $k=1,2,3$ & exact\\
\bottomrule
\end{tabular}}
\medskip
\caption{Verification (\texttt{cohn\_localization.py}) on the $K_4$ Iwahori tower and the
$(1,1,1)$ abelian tower.}
\label{XIII:tab:battery}
\end{table}

\section{Scope, residue, corrections}\label{XIII:sec:scope}

\begin{remark}[What a main conjecture would be]\label{XIII:rem:main}
In the theory of \cite{CFKSV} a main conjecture asserts that an \emph{analytic} object --- a
$p$-adic $L$-function characterized by interpolation of special values, constructed
independently of the module --- equals the algebraic characteristic element, an element of
$K_1$ of the localized Iwasawa algebra whose boundary is the class of the Selmer module. The
algebraic half is what this paper constructs: $[\Theta_{\cH}]\in K_1(\HS)$ with
$\partial[\Theta_{\cH}]$ the class of the tower's cokernels, the analogue of ``$\partial\xi=[X]$''.
The analytic half does not exist yet: there is no independently defined element of
$K_1(\HS)$ interpolating the $L$-values of the tower --- the candidates are the pseudo-determinants
$n_k\kappa_k=\prod_{j}(q{+}1-\lambda_j)$, products of local $L$-values at the Eisenstein point
as in [Ch.~\ref{chap:I}, Thm.~3.5], but a construction of a single element from them, rather than from
$\Delta$, is exactly what is missing. The non-commutative main conjecture for Hecke towers
is therefore: \emph{construct $\mathcal L_{\cH}\in K_1(\HS)$ from the $L$-values alone and
prove $\mathcal L_{\cH}=[\Theta_{\cH}]$ modulo the image of $K_1(\cH)$}. Its evaluated form
is [Ch.~\ref{chap:VII}, Thm.~2.1]; its $K_1$ form is open.
\end{remark}

\begin{remark}[Solved, open, impossible]\label{XIII:rem:scope}
\emph{Solved:} the residue of [Ch.~\ref{chap:V}, Prob.~5.1] and [Ch.~\ref{chap:VII}, Rem.~3.1] in the form asked --- a single
class in $K_1$ of a localized degeneracy algebra whose level evaluations and boundary
reproduce the snake factorization, with (A1) and (A2) separated as evaluation and Eisenstein
subclass and (A3) as the abelian reduction. \emph{Open:} (a) the analytic side
(Remark~\ref{XIII:rem:main}); (b) stable flatness of $\cH\to\HS$ and hence the higher localization
sequence \cite{NR}; (c) the structure of $\HS$ itself --- whether it is an order in a
semisimple algebra or has torsion, and the image of $K_1(\cH)$ in $K_1(\HS)$, which
measures how far $[\Theta_{\cH}]$ is from unique; (d) the rank-two version over the
degeneracy algebra of [Ch.~\ref{chap:XII}], where no abelian sector exists. \emph{Impossible or empty as
stated:} ``$\partial[\Theta]=[X]$'' as a theorem with content (Theorem~\ref{XIII:thm:tauto}).
\end{remark}

\begin{remark}[Corrections to the task specification]\label{XIII:rem:corrections}
(1) The identity $\partial([\Theta_{\cH}])=[X_\infty]$ is the definition of the boundary map in
the relative sequence, not a theorem; the paper proves what does have content
(Theorems~\ref{XIII:thm:eval}--\ref{XIII:thm:abelian}). (2) The class must be formed on the
augmentation-zero lattices: on $\Zp^{n_k}$ the Laplacian has the Eisenstein kernel and is not
$\Sigma$-inverting; the specification's ``$\Delta_\infty\in M_n(\widehat{\cH}_\infty)$'' has no
$n\times n$ presentation because $n_k$ grows (contrast [Ch.~\ref{chap:IV}], where $\Zl[G_k]^{n}$ is free of
constant rank). (3) The Euler term is not produced by ``the connecting map'' inside $K_1$;
it is the Eisenstein subclass $[\Z/n_k]$ of the boundary, and it is invisible to the
$K_1$-evaluations, which see the arc term only. (4) ``$\Sigma$ = the set of morphisms
inducing surjections with finite cokernel on $X_\infty$'' is replaced by the single element
$\Delta_{(K)}$; universal localization at one element already gives the class. (5) The
higher $K$-theory sequence needs stable flatness \cite{NR}; only the low-degree sequence
\cite{Sch} is used, and only its first map matters for the construction.
\end{remark}

\section*{Acknowledgement of provenance and caveats}

Unrefereed preprint. Every numbered statement is proved from Papers III--XII and finite
linear algebra, given the existence of universal localization \cite{Co} and the low-degree
exact sequence \cite[Ch.~4]{Sch}, both quoted at chapter level (the precise theorem numbers
were not consulted). Bibliographic data of \cite{Co,Sch,NR} were checked against publisher
or indexing pages; \cite{CFKSV} is the entry verified in [Ch.~\ref{chap:V}].
